\section{Introduction}

Web-agent benchmarks are becoming more realistic, but most still treat website change as a fixed nuisance: choose a set of versions or perturbations, run an agent, and report success rate. That protocol measures brittleness, yet it does not answer the design question that appears after each evaluation. Once an agent has a strong rulebook, which website changes still expose its capability boundary, and how should the next rule update be selected?

This question matters for reflection-based agents. A reflection rule such as ``preserve a login redirect'' or ``stop when the target state is already visible'' only helps if the benchmark contains valid situations where that behavior is required. If generated websites are too easy, they provide no learning signal. If they are broken, failures cannot support a rule. \webcoevo{} studies the middle ground: use the current agent as a solver, generate websites that target its observed frontier, and mine small cross-version reflection rules from paired trajectory evidence.

We instantiate this loop on \linkding{}, a self-hosted bookmark manager with deterministic local deployment and code-level website overlays. The agent is Qwen3-VL in a BrowserGym/AgentLab/UI-TARS stack. The base policy can receive \expel{} task-experience rules distilled from earlier trajectories, while the cross-version layer injects compact reflection rules about how to adapt when the website changes. This separation is important. \expel{} rules tell the agent what helped on related tasks; cross-version reflection rules tell it how to preserve task intent across a shifted interface.

The upgraded WebCoEvo evidence has three website versions. The first is \linkding{} V1.45.0, the original release. The second is a moderately modified website V2 that tests whether R1 reflection rules improve over \expel{}-style prompting under moderate website drift. The third is a harder generated website suite V3 with seven drift categories: \texttt{access}, \texttt{surface}, \texttt{content}, \texttt{runtime}, \texttt{process}, \texttt{structural}, and \texttt{functional}. Website V3 is the primary stress test because \expel{}-style prompting collapses there even when it performs well on \linkding{} V1.45.0 and website V2.

The main empirical result is sharper than in the earlier draft. On the 68 training tasks from \linkding{} V1.45.0, \expel{}-style rules raise success from \(9/68=13.2\%\) to \(56/68=82.4\%\). On website V2, R1 reflection rules improve over \expel{}-style prompting from \(60/68=88.2\%\) to \(67/68=98.5\%\) on the training tasks and from \(114/167=68.3\%\) to \(143/167=85.6\%\) on the held-out validation tasks. On website V3, however, \expel{}-style prompting drops to \(8/68=11.8\%\) on the training tasks and \(66/167=39.5\%\) on the held-out validation tasks. R1 recovers most of the training split, reaching \(65/68=95.6\%\), and transfers to the held-out validation tasks at \(97/167=58.1\%\). R4, a GPT-5.4-assisted update built from transition evidence, ties R1 on the held-out validation tasks but reaches only \(60/68=88.2\%\) on the training tasks, so iterative reflection is useful but not monotonic.

This paper makes three contributions.
\begin{enumerate}
\item We formulate a capability-targeted adversarial co-evolution protocol for web-agent reflection rules: profile current failures, generate drifted websites at the boundary, audit task validity, mine behavior gaps, and promote only compact rule updates that survive validation.
\item We materialize the protocol in \webcoevo{}, a standalone \linkding{} runner with explicit task normalization, reset logic, prompt construction, rulebook selection, and legacy eval/trace JSONL export.
\item We report the upgraded evidence across \linkding{} V1.45.0, website V2, and website V3: \expel{}-style rules help on the original site, but compact reflection rules are what preserve performance under generated website drift, with R1 strongest on the training tasks under website V3 and R1/R4 tied on the held-out validation tasks under website V3.
\end{enumerate}

The claim boundary is deliberately narrow. We do not claim universal web-agent robustness, nor that later reflection-rule versions always improve. The evidence is single-site, and several drift families remain difficult on the held-out validation tasks under website V3. The contribution is an auditable loop that turns those failures into the next environment and rule-design targets.
